\documentclass[11pt,a4paper]{article}

\usepackage[margin=1in]{geometry}
\usepackage{amsmath,amssymb}
\usepackage{booktabs}
\usepackage{longtable}
\usepackage{array}
\usepackage[protrusion=true,expansion=false]{microtype}
\usepackage{xcolor}
\usepackage{titlesec}
\usepackage{enumitem}
\usepackage{fancyhdr}
\usepackage[breakable]{tcolorbox}
\usepackage[numbers,sort&compress]{natbib}
\usepackage[hidelinks,colorlinks=true,linkcolor=linknavy,urlcolor=linknavy,citecolor=linknavy]{hyperref}

\definecolor{linknavy}{HTML}{184A7A}
\definecolor{rulegray}{HTML}{9A9A94}
\definecolor{boxbg}{HTML}{F4F2EC}
\definecolor{warnbg}{HTML}{FBF0E8}

\tcbset{
  keybox/.style={colback=boxbg, colframe=rulegray, boxrule=0.4pt, arc=2pt,
    left=8pt, right=8pt, top=6pt, bottom=6pt, breakable},
  warnbox/.style={colback=warnbg, colframe=orange!60!black, boxrule=0.6pt, arc=2pt,
    left=8pt, right=8pt, top=6pt, bottom=6pt, breakable}
}

\titleformat{\section}{\large\bfseries}{\thesection}{0.6em}{}
\titleformat{\subsection}{\normalsize\bfseries}{\thesubsection}{0.6em}{}
\titlespacing*{\section}{0pt}{1.7ex plus .4ex}{0.9ex}

\setlength{\parskip}{0.55em}
\setlength{\parindent}{0pt}
\renewcommand{\arraystretch}{1.2}

\pagestyle{fancy}
\fancyhf{}
\fancyhead[L]{\small\itshape Methodology of an AI-assisted identifiability audit}
\fancyhead[R]{\small\thepage}
\renewcommand{\headrulewidth}{0.2pt}

\title{\vspace{-1.5em}\textbf{Seven Refutations and a Missing Host}\\[0.4em]
\large A Methodological Post-Mortem of an AI-Assisted Research Project\\
\normalsize\itshape What the process caught, what it could not have caught,\\
\itshape and what that implies for human and machine research practice}

\author{Eric Schultz\\[0.2em]
\small Independent Researcher --- ORCID 0009-0006-6283-1696\\[0.5em]
\small\itshape Prepared with a five-perspective review panel:\\
\small\itshape Research Methodologist, Cognitive Scientist, Reproducibility Auditor,\\
\small\itshape Red-Team Analyst, and Document Architect}

\date{\small Companion methodology paper to \emph{An Identifiability Audit of Evolutionary Oncology}}

\begin{document}
\maketitle
\thispagestyle{fancy}

\begin{abstract}
\noindent
We document the complete methodology of a research project conducted in close
collaboration with a large language model, in which an identifiability audit of
evolutionary oncology was designed, executed, refuted, re-derived, and refuted again
across a single extended session. The project produced a substantive scientific result
--- that the parameters underpinning adaptive cancer therapy are unidentifiable under
standard-of-care measurement, and that the field has been arguing about the wrong
parameter --- but the process is at least as instructive as the product.

\textbf{Seven distinct hypotheses were refuted, six of them by the authors' own
adversarial self-testing.} Every error was one of \emph{interpretation}, never of
\emph{implementation}: an independent verification suite found the code correct to
$10^{-6}$ throughout. This asymmetry is the paper's central empirical claim about
method, and it inverts where most research effort is currently spent.

\textbf{The seventh refutation was found by an external human reviewer, and could not
have been found internally.} The model had omitted the host organism from a model of
drug holidays in cancer patients. No amount of internal verification --- multi-start
optimisation, Fisher-information cross-checks, independent reimplementation, null
controls --- can detect a missing compartment, because all of it operates \emph{inside}
the specification under audit. We treat this as the load-bearing finding of the
methodology and argue it generalises: \textbf{AI-assisted research systematically
converts interpretation errors into detectable ones while leaving omission errors
undetectable, and therefore concentrates residual risk in exactly the category humans
are worst at auditing.}

We provide a comparative analysis against standard human research practice, an
adversarial threat model of the AI-assisted pipeline (including a documented,
self-caught instance of citation fabrication), a reproducibility statement, and a set
of concrete protocols we recommend to any researcher adopting this mode of work.
\end{abstract}

\tableofcontents

\vspace{1em}
\hrule
\vspace{0.4em}

\section{Introduction and scope}

This is a methodology paper, not a results paper. The scientific object it refers to
--- an identifiability audit of adaptive cancer therapy --- is described in a companion
manuscript \citep{schultz2026audit} and summarised here only where the process cannot
be understood without it.

The claim we defend is narrow and, we think, testable: that \textbf{an AI-assisted
research loop has a characteristic and predictable error profile}, distinct from that
of an unassisted human researcher, and that this profile can be described precisely
enough to be defended against. It is not the profile the popular discussion expects.
The failures were not hallucinated facts or incoherent mathematics. The mathematics was
correct every single time. The failures were failures of \emph{framing} --- of asking
the wrong question with impeccable rigour --- and, in one decisive case, of omitting a
piece of the world from the model entirely.

We report the failures in full, in chronological order, including one in which the model
fabricated bibliographic metadata and caught itself doing so. A methodology paper that
reported only the successes of the method would be committing the precise error the
method exists to prevent.

\vspace{0.5em}
\hrule
\vspace{0.6em}
\begin{center}\large\bfseries PART I --- THE RECORD\end{center}
\vspace{0.2em}
\hrule

\section{Chronology of the research process}
\emph{Contributed by the Research Methodologist \& Historian.}

\subsection{Phase 0: an open question, deliberately underspecified}

The project did not begin with a hypothesis. It began with an open prompt --- what is
structurally underemphasised in cancer research? --- and a list of candidate deficits:
evolutionary dynamics ignored in favour of static targeting; metastasis and dormancy
understudied relative to primary tumour biology; preclinical irreproducibility
\citep{errington2021investigating,begley2012raise}; prevention underfunded relative to
treatment; overdiagnosis; the absence of any institution funded to run unprofitable
trials; and a coarse measurement cadence.

The choice to pursue the last of these was made on tractability grounds, not importance
grounds. This is worth recording because it is a species of selection effect: \textbf{the
question we could attack with the tools at hand became the question we attacked.} We
have no argument that it was the most important item on the list.

\subsection{Phase 1--7: the hypothesis chain}

Table~\ref{tab:chain} is the complete record. Each row is a hypothesis that was
seriously held, operationalised in code, and abandoned.

\begin{table}[h]
\centering
\small
\begin{tabular}{p{0.9cm}p{5.4cm}p{7.4cm}}
\toprule
& \textbf{Hypothesis held} & \textbf{What refuted it}\\
\midrule
H1 & Oncology samples too slowly; denser ctDNA sampling closes the gap. \emph{(The founding thesis.)} & Direct test. 6 draws $\to$ 281 draws (daily, 40 weeks): sign-error rate $57\% \to 57\%$, \emph{exactly unchanged}. The system was not bandwidth-limited but information-limited.\\
\addlinespace
H2 & Clone number is recoverable with a per-variant panel; the polyclonal hazard is closable. & The candidate set. We had fitted $k=1,2,3,4$ --- all discrete. Admitting a \emph{continuum} model, BIC preferred it in \emph{both} worlds, including 6/8 against a genuinely 3-clone truth.\\
\addlinespace
H3 & The fitness cost of resistance $c$ is the quantity to measure. \emph{(The field's framing, inherited unaudited.)} & A sweep we had no particular reason to run. With $c = 0$ --- resistance entirely free --- adaptive therapy still beat MTD substantially. The cost was not the mechanism; competition ($\alpha$) was.\\
\addlinespace
H4 & The benefit of adaptive therapy is flat in $c$; the fitness cost is irrelevant. & Our own arithmetic. We had normalised by MTD survival, which \emph{also} rises with $c$. In absolute days the benefit roughly doubles. A moving denominator concealed a real effect.\\
\addlinespace
H5 & The directional suppression coefficient $\alpha_{SR}$ is structurally unidentifiable under any design. & Two errors. (i) We applied a \emph{per-patient} standard to $\alpha_{SR}$ that no other parameter had been held to. Pooled at $n=10$ it resolves cleanly. (ii) The bound came from pseudo-inverting a Fisher matrix with $\kappa \sim 10^7$, and was \emph{non-monotone in the data} --- a numerical red flag we read as a finding.\\
\addlinespace
H6 & Resource supply rate can be used to vary competition strength in a mechanistic model. & The $R^*$ rule \citep{tilman1982resource}. Equilibrium resource level is independent of supply. There was never a no-competition control. The mis-specified experiment produced a clean-looking, entirely spurious $11\times$ inversion of the headline result.\\
\addlinespace
\textbf{H7} & \textbf{Any incidental treatment interruption is a free excitation of the system.} & \textbf{An external human reviewer.} A Grade 3/4 toxicity hold is not a clean drug-off window --- it is a systemic crisis that alters the very growth and competition parameters being measured. \textbf{No internal method found this, and none could have.}\\
\bottomrule
\end{tabular}
\caption{The complete hypothesis chain. Six refutations were internal; the seventh, and
most consequential, was not.}
\label{tab:chain}
\end{table}

\subsection{The three genuine pivots}

Three moments changed the direction of the work rather than merely correcting it.

\textbf{(1) Reframing from biology to information.} The decisive move was to stop asking
\emph{does a fitness cost exist?} --- a biological question the field has debated for a
decade --- and instead ask \emph{could we tell, from the data we collect?} This is a
question about Fisher information, not about cancer, and it is answerable in an
afternoon. The reframing is the entire reason the project produced anything.

\textbf{(2) The $R^*$ rule.} Realising that equilibrium resource concentration is
independent of supply rate \citep{tilman1982resource} invalidated an entire mechanistic
experiment \emph{that had already produced a confident, publication-shaped table}. The
question that caught it --- ``what does this knob actually \emph{do} to the
equilibrium?'' --- should have been asked before running the sweep, not after.

\textbf{(3) Independent corroboration we did not expect.} Only during the citation pass
did we discover \citet{bacevic2017spatial}, who report experimentally that
\emph{spatial competition alone} constrains resistance to targeted therapy. Our audit
had derived from scratch that competition, not fitness cost, drives adaptive therapy.
An existing experiment already said so, and the theoretical literature had not
connected it. This converted a surprising simulation result into a simulation result
that agrees with a real experiment.

\vspace{0.5em}
\hrule
\vspace{0.6em}
\begin{center}\large\bfseries PART II --- ANALYSIS\end{center}
\vspace{0.2em}
\hrule

\section{Comparative analysis: the AI loop versus standard practice}
\emph{Contributed by the AI Safety \& Cognitive Scientist.}

\subsection{The economics of refutation}

The single largest methodological difference is not intelligence. It is
\textbf{the marginal cost of attempting to destroy your own claim.}

In standard practice, testing H1 --- does denser sampling actually help? --- requires
writing a simulation, validating it, running it, and analysing it: days to weeks of a
researcher's time, and the researcher already believes the answer. The rational move,
under realistic incentives, is not to check. In the loop described here, the same test
cost roughly four minutes and no emotional investment. \textbf{When refutation becomes
cheap, refutation becomes rational} --- and a thesis that would have survived to
publication instead died in the first hour.

This effect compounds. Six internal refutations occurred in a single session. No
individual step was beyond a competent human; the difference is that a human would have
had to \emph{want} to take all six, in sequence, against their own accumulating
investment.

\begin{tcolorbox}[keybox]
\textbf{The core asymmetry.} Every error in this project was one of
\emph{interpretation}. Not one was a coding error. An independent verification suite
--- three model implementations agreeing to $<10^{-6}$; the piecewise integrator against
brute-force Euler to $1.3\times10^{-4}$; the Fisher information against the numerical
Hessian of the expected log-likelihood to $2\times10^{-8}$ --- passed cleanly at every
stage, \textbf{including while the conclusions it was supporting were wrong}.

Research culture invests heavily in code review, unit tests, and computational
reproducibility. Those defences were, in this project, \textbf{100\% effective and
0\% relevant}.
\end{tcolorbox}

\subsection{What human researchers should actually adopt}

We offer six protocols. They are not AI-specific; the AI merely made them cheap enough
to notice they were missing.

\begin{enumerate}[leftmargin=1.8em,itemsep=0.4em]

\item \textbf{Run the identifiability audit \emph{before} the experiment, not after the
disagreement.} The entire contribution of the companion paper is that a decade of
argument was conducted over parameters that, under the field's own model and its own
data, could not have been measured either way. Ask what your protocol can \emph{see}
before you argue about what your model \emph{implies}
\citep{raue2009structural,eisenberg2017confidence}.

\item \textbf{Audit the candidate set, not just the model.} H2 died because our model
selection could only choose among the models we offered it. Concluding ``the tumour is
discrete'' from a menu of exclusively discrete models is circular. \textbf{Before
running model selection, write down the class of model you have made it impossible to
select.}

\item \textbf{Treat every normalisation as a suspect.} H4 --- reporting a
\emph{relative} gain whose denominator was itself moving with the parameter of interest
--- is the most ordinary statistical error in the record, and it produced the most
confident sentence in the draft it appeared in. This is a special case of the garden of
forking paths \citep{gelman2013garden}: the normalisation was not chosen dishonestly,
it was chosen \emph{because it looked clean}, immediately after we became excited.

\item \textbf{Distrust bounds computed on ill-conditioned surfaces --- and check
monotonicity as a diagnostic.} H5's spurious standard errors were \emph{non-monotone in
the quantity of data}. More data cannot reduce information. That non-monotonicity was
the numerics announcing that the bound was meaningless, and we read it as a substantive
finding. \textbf{A cheap, general check: if your uncertainty does not shrink when you
add data, your uncertainty estimate is broken, not interesting.}

\item \textbf{Ask what every knob controls before you turn it.} H6 produced a
fully-formed spurious result because we varied a parameter that, by the $R^*$ rule,
does not control the thing we believed it controlled. Nothing in the output looked
wrong.

\item \textbf{Maintain an errata log as a first-class artefact.} The companion paper
carries its seven refutations in the body, not a supplement. This is unusual and
strategically risky. It is also the only reason the eighth claim is worth anything: a
reader can see exactly how many times we were wrong, and by what mechanism.

\end{enumerate}

\subsection{What the loop is bad at, and it is not what you expect}

The AI-assisted loop was strong at exactly the things folklore says it is weak at:
arithmetic, code correctness, internal consistency, and the willingness to abandon a
cherished claim (there being no cherishing). It was weak at exactly one thing, and that
weakness is the subject of the next two sections: \textbf{it could not notice what was
not there.}

\section{Adversarial stress-testing and threat modelling}
\emph{Contributed by the AI Red Teaming \& Vulnerability Analyst.}

We evaluate the pipeline as an attacker would: not asking whether it produced a good
result, but asking \textbf{what class of wrong result it would produce with maximum
confidence and minimum detectability.}

\subsection{Threat model}

\begin{table}[h]
\centering
\small
\begin{tabular}{p{3.5cm}p{4.6cm}p{5.6cm}}
\toprule
\textbf{Failure mode} & \textbf{Did it occur?} & \textbf{Detection / mitigation}\\
\midrule
\textbf{Confirmation-bias feedback loop} & \textbf{Yes --- H4.} A normalisation was chosen that flattered a conclusion reached minutes earlier. & Caught internally, but only on a later unrelated audit. \emph{Mitigation:} pre-register the estimand and the normalisation before running the sweep.\\
\addlinespace
\textbf{Fabrication / hallucination} & \textbf{Yes --- bibliography.} Author lists truncated in sources (``Diehl F, Schmidt K, Choti MA, \emph{et al.}'') were silently completed from parametric memory across 7 entries. One reconstruction propagated a name error. & Caught by an explicit verification pass in which every author name was required to trace to a source actually read. All 7 entries were truncated to verified authors $+$ \texttt{and others}. \emph{This is the canonical LLM failure mode and it occurred exactly where attention was lowest.}\\
\addlinespace
\textbf{Fluency-as-evidence} & \textbf{Yes --- H6.} The mis-specified mechanistic model produced a clean, well-formatted, publication-shaped table pointing the \emph{opposite} direction from the truth. & \textbf{Nothing in the output looked wrong.} No internal check flags a well-executed experiment on the wrong system. \emph{Mitigation:} require a mechanistic account of what each swept parameter controls at equilibrium.\\
\addlinespace
\textbf{Sycophancy / drift toward the user's framing} & \textbf{Structurally present.} The model inherited the field's cost-of-resistance framing (H3) uncritically, and would plausibly have inherited the user's had one been asserted. & Partially mitigated by the user posing open rather than leading questions. \emph{Unresolved risk.}\\
\addlinespace
\textbf{Monoculture / correlated blind spots} & \textbf{Yes --- H7, decisively.} The same system generated the hypothesis, the model, the code, \emph{and} the critique. Its blind spots are therefore perfectly correlated across all four roles. & \textbf{Not mitigable from inside.} See \S\ref{sec:h7}.\\
\addlinespace
\textbf{Prompt injection via retrieved content} & Not observed. Literature search returned no adversarial content. & Live vector in any pipeline with automated retrieval. Treat fetched text as untrusted data, never as instruction.\\
\addlinespace
\textbf{Optimiser-assisted confirmation} & Guarded against. & Multi-start MLE with restarts drawn from the \emph{entire} admissible box, never perturbed from the truth; $c$ free on $[-3, 0.99]$ so a fit was \emph{permitted} to conclude resistance is beneficial.\\
\bottomrule
\end{tabular}
\caption{Threat model. Five of seven anticipated failure modes actually occurred.}
\end{table}

\subsection{Red-team protocols that were executed}

\begin{itemize}[leftmargin=1.4em,itemsep=0.25em]
\item \textbf{Independent reimplementation.} Three separate implementations of the
dynamics, cross-checked to $<10^{-6}$.
\item \textbf{Analytic-versus-numerical cross-check.} Fisher information matrix versus
the numerical Hessian of the expected negative log-likelihood: agreement to
$2\times10^{-8}$.
\item \textbf{Integrator validation.} Event-driven piecewise LSODA against a
brute-force 20{,}000-step Euler integration: $1.3\times10^{-4}$.
\item \textbf{Observable saturation check.} Verified the allele-fraction channel never
pinned at its floor or ceiling, which would have silently destroyed sensitivity.
\item \textbf{Seed-set stability.} Headline rates recomputed across independent seed
sets; this demoted several point estimates to ranges (50--57\%, not 57\%) and caught one
near-miss where a 57\% and a 14\% figure were both drawn from $n=14$ replicates.
\item \textbf{A genuine null control.} The MacArthur model's $\omega = 0$ arm --- clones
consuming \emph{different} resources --- is a real zero-competition control, and returns
exactly zero benefit. Its absence was precisely what invalidated H6.
\end{itemize}

\subsection{The protocol that was missing, and had to come from outside}
\label{sec:h7}

\begin{tcolorbox}[warnbox]
\textbf{Every protocol above operates \emph{inside} the model specification.}

Multi-start optimisation searches the parameter space of a given model. Cross-checking
the Fisher information against the Hessian verifies that two derivations of the
\emph{same} likelihood agree. Independent reimplementation confirms that three encodings
of the \emph{same} dynamics coincide. A null control tests a hypothesis
\emph{within} the specified system.

\textbf{None of them can detect that the system is missing a compartment.} They will all
pass, cleanly and to eight decimal places, on a model of a tumour that has no patient
attached to it.
\end{tcolorbox}

This is what happened. The design's central move is to exploit incidental treatment
interruptions as ``free excitations'' of the tumour system. An external reviewer with
clinical knowledge observed that a hold for Grade 3/4 toxicity is not a clean
perturbation --- the patient is in systemic crisis (hepatotoxicity, neutropenia,
inflammation), and the growth and competition parameters of a tumour in a metabolically
deranged host are not those of a tumour in a well one.

When we modelled this, the critique was confirmed and the structure of the damage proved
sharper than the critique itself. A \emph{symmetric} host insult is comparatively benign
--- the fitness cost is a \emph{ratio}, and a common factor largely cancels. A
\emph{differential} insult that spares the resistant clone is fatal: the sign-error rate
returns to \textbf{58\%}, and, worst of all, the bias runs \textbf{downward} --- toward
concluding that resistance is free, which is \emph{precisely} the error the entire paper
exists to prevent. The confound is aligned with the failure mode. Fitting the host effect
as a nuisance parameter does not rescue it ($36\% \to 22\%$).

The remedy is an inclusion criterion, not a statistical one: only interruptions that
perturb the drug \emph{and nothing else} --- insurance lapse, supply interruption,
scheduled surgery, non-adherence --- are clean excitations.

\begin{tcolorbox}[keybox]
\textbf{The generalisable claim.} An AI-assisted loop with adequate verification
converts \emph{interpretation errors} into \emph{detectable} errors at very low cost.
It does nothing whatsoever about \emph{omission errors} --- and by dramatically raising
the polish, internal consistency, and apparent rigour of the artefact, it makes
omission errors \emph{harder} for downstream readers to detect.

\textbf{Residual risk is therefore concentrated entirely in the category the method
cannot see, and delivered in the most persuasive possible packaging.}
\end{tcolorbox}

\subsection{Recommended protocols for future work}

\begin{enumerate}[leftmargin=1.8em,itemsep=0.3em]
\item \textbf{Mandatory external domain review, treated as a required experiment rather
than a courtesy.} H7 was worth more than any internal check performed. Budget for it
accordingly.
\item \textbf{Heterogeneous adversaries.} A critique generated by the same model that
generated the work inherits its blind spots exactly. Use a different architecture, a
different training distribution, or --- best --- a human with different training.
\item \textbf{An explicit ontology review.} Before verification begins, ask: \emph{what
entities exist in the real system that do not appear in this model?} For this project
the correct answer included ``the patient'', and no one asked.
\item \textbf{Pre-register the estimand and the normalisation.} Defuses H4-class errors.
\item \textbf{Require every cited fact to trace to a source actually retrieved in-session.}
Defuses the bibliography failure. Parametric memory is not a citation.
\end{enumerate}

\section{Limitations and broader impacts}
\emph{Contributed by the Ethics \& Reproducibility Auditor.}

\subsection{Limitations of the scientific object}

The companion paper is \textbf{entirely simulation-based}. It makes no empirical claim
about any patient, and its central assertion --- that certain parameters are
unidentifiable under current measurement --- is a statement about information in a model,
not a discovery about cancer. The efficiency figures ($n \approx 7$--$10$ patients) are
\emph{ideal-pooling bounds} and will inflate under realistic between-patient
heterogeneity; the qualitative claim (tens, not tens of thousands) is robust, the integer
is not, and it is the integer a reader will quote.

\subsection{Limitations and hazards of the methodology}

\begin{itemize}[leftmargin=1.4em,itemsep=0.35em]

\item \textbf{The method raises the production value of error.} This is the principal
societal hazard, and it deserves to be stated plainly. A wrong result that is
well-typeset, correctly cited, internally consistent, and accompanied by a passing
verification suite is \emph{more} dangerous than a wrong result that looks wrong. H6
produced exactly such an artefact. Peer review's informal defences --- sloppiness as a
signal of unreliability --- are substantially degraded by this class of tool.

\item \textbf{Automation bias in reviewers.} A manuscript displaying this density of
self-refutation and verification may induce unearned trust. \textbf{The errata section is
not evidence that the remaining claims are correct.} It is evidence only that a
particular class of error was searched for.

\item \textbf{Speed outruns domain contact.} The project reached a fully-formed,
cited, typeset manuscript before a single domain expert had seen it --- which is
precisely how a model of cancer patients came to have no patients in it. The traditional
slowness of research is partly incidental friction and partly a \emph{sampling process
over human expertise}. Removing the friction removes the sampling.

\item \textbf{Selection effect on the research question.} We chose the deficit we could
attack, not the deficit that mattered most (\S2.1). AI assistance sharpens this bias by
making formally tractable questions dramatically cheaper relative to empirically
grounded ones.

\item \textbf{Documented fabrication incident.} During bibliography construction, author
lists that appeared truncated in retrieved sources were completed from parametric memory
in seven entries; one reconstruction propagated a name error. This was caught by an
explicit verification pass and corrected. \textbf{We report it because a methodology
paper that concealed it would be worthless.} Any reader adopting this workflow should
assume this failure will occur in their pipeline, because it will.

\item \textbf{Single-operator, single-session.} $n = 1$. Everything here is a case study.
We make no claim that these findings replicate across researchers, models, or domains,
and we would regard a confident generalisation from this sample as an instance of the
error the paper documents.
\end{itemize}

\subsection{Ethical considerations}

No human subjects, no patient data, and no personally identifiable information were
involved at any stage; all data were synthetic. The proposed clinical design is
\textbf{purely observational by construction} --- it prescribes no change to any
treatment decision --- which was an explicit design constraint from the outset rather
than a post-hoc justification. Nevertheless, we note that a compelling, well-produced
simulation argument could influence clinical research priorities, and that this
influence would be unearned: \textbf{no result in the companion paper constitutes
evidence about a patient.}

\section{Reproducibility statement}

\begin{table}[h]
\centering
\small
\begin{tabular}{p{4.3cm}p{9.5cm}}
\toprule
\textbf{Artefact} & \textbf{Location / value}\\
\midrule
Code repository & \texttt{[PLACEHOLDER: repository URL]}\\
Archived snapshot & \texttt{[PLACEHOLDER: Zenodo DOI]}\\
Companion manuscript & \texttt{[PLACEHOLDER: preprint DOI]}\\
Simulation scripts & 27 Python modules: dynamics, Fisher information, bootstrap MLE, cohort pooling, ontology selection, scheduling, mechanistic replication, toxicity confound, verification suite\\
Data & \textbf{None.} All results are generated from synthetic data; no external dataset is used or required.\\
Random seeds & Explicitly set and reported per experiment; headline results verified across independent seed sets\\
Language / libraries & Python 3; NumPy, SciPy (\texttt{solve\_ivp} LSODA, \texttt{least\_squares}), Matplotlib\\
Hardware & Single CPU core; no GPU\\
Total compute & $<1$ CPU-hour. \emph{Energy cost of the simulations is negligible; the dominant environmental cost of this project is LLM inference, which we are unable to quantify from available disclosures.}\\
Verification suite & \texttt{verify.py} --- cross-implementation, integrator, Fisher-vs-Hessian, saturation, seed-stability. All checks pass.\\
Session transcript & \texttt{[PLACEHOLDER: link, if released]} --- \emph{we recommend release; the errors are the contribution}\\
\bottomrule
\end{tabular}
\caption{Reproducibility statement. Placeholders to be completed at deposition.}
\end{table}

\textbf{Reproduction instructions.} Every numerical claim in the companion paper is
regenerated by executing the corresponding module. No result depends on proprietary
data, licensed software, or specialised hardware. The verification suite should be run
first; if it fails, no downstream number should be trusted.

\textbf{A caveat specific to this workflow.} Reproducibility of the \emph{code} is not
reproducibility of the \emph{research}. Every error in this project would reproduce
perfectly. A reader re-running our scripts would obtain our wrong answers to eight
decimal places, seven times over. \textbf{Computational reproducibility is orthogonal to
correctness, and this project is an unusually clean demonstration of that fact.}

\section{Conclusion}

The scientific result of this project is that evolutionary oncology cannot currently
measure the parameters its models depend on, and has been arguing about the wrong one.
The methodological result is more uncomfortable.

An AI-assisted loop refuted six of its own hypotheses in a single session --- a rate no
unassisted researcher would sustain, not through lack of ability but through entirely
rational unwillingness to spend weeks disproving their own work. It maintained perfect
code correctness throughout, verified to eight decimal places. It produced a polished,
cited, internally consistent manuscript.

And it left the patient out of a model of patients.

\begin{tcolorbox}[keybox]
\textbf{Internal rigour cannot find a missing compartment.} Every verification protocol
we ran --- and we ran good ones --- operates inside the specification under audit. They
were 100\% effective and 0\% relevant to the one error that could have invalidated the
work.

The correct lesson is not that AI-assisted research is unreliable. It is that its
reliability profile is \textbf{sharply non-uniform}: near-perfect where verification
applies, and \emph{unchanged} where it does not --- while the persuasiveness of the
output rises in both cases.

The defence is not more verification. It is a human who knows what the model left out.
\end{tcolorbox}

\vspace{0.5em}

We would encourage anyone adopting this mode of work to keep an errata log, publish it,
and treat the first domain expert who reads their manuscript as the most important
experiment they will run.

\vspace{1em}
\hrule
\vspace{0.6em}

{\small\itshape Correspondence: ORCID 0009-0006-6283-1696. The authors declare no
competing interests. Entries marked \texttt{[PLACEHOLDER]} are to be completed at
deposition. Bibliographic entries marked \texttt{[UNVERIFIED]} in the source
\texttt{.bib} have not been checked against primary sources in-session and must be
verified before submission --- an instruction we issue in light of \S5.2.}

\bibliographystyle{unsrtnat}
\bibliography{methodology}

\end{document}
